\subsection{分圆多项式的不可约性}

设 p 为 K 的特征指数，且 n 为与 p 互素的整数。我们用 \(\Phi_{n} \in \mathbf{K}[\mathbf{X}]\) 表示整系数多项式 \(\Phi_{n}\) 在将 X 映到 X 的、从 Z{[}X{]} 到 K{[}X{]} 的唯一同态下的像。

引理3. --- \(\Phi_n\) 在 \(\mathbf{K}_s\) 中的根是本原的 \(n\) 次单位根。

用 \(S_{n}\) 表示 \(\Phi_{n}\) 在 \(K_{s}\) 中的根组成的集合。根据公式（6），集合 \(\mu_{n}(K_{s})\) 是诸 \(S_{d}\) （对整除 n 的 d）的并。于是每个本原 n 次单位根都属于 \(S_{n}\) ，引理现在由命题4（V，第81页）得出。

命题6. --- 设 p 为 K 的特征指数，且设 \(n \geqslant 1\) 为与 p 互素的整数。多项式 \(\Phi_{n}(\mathbf{X})\) 在 K{[}X{]} 中不可约，当且仅当同态 \(\chi_{n}: \operatorname{Gal}(K_{s}/K) \to (\mathbf{Z}/n\mathbf{Z})^{*}\) 是满射的。

根据引理3，我们有 \(\mathbf{R}_n(\mathbf{K}) = \mathbf{K}(\zeta)\) ，对每个根 \(\zeta\) （ \(\Phi_n(\mathbf{X})\) 的根）而言，因此 \(\Phi_n(\mathbf{X})\) 在 \(\mathbf{K}[\mathbf{X}]\) 中不可约，当且仅当次数 \(\varphi(n)\) （ \(\Phi_n(\mathbf{X})\) 的）等于 \([\mathbf{R}_n(\mathbf{K}) : \mathbf{K}]\) 。进一步地， \(\mathbf{R}_n(\mathbf{K})\) 在 \(\mathbf{K}\) 上的伽罗瓦群的阶为 \([\mathbf{R}_n(\mathbf{K}) : \mathbf{K}]\) ，且它同构于 \((\mathbf{Z}/n\mathbf{Z})^*\) 的一个子群，即 \(\chi_n\) 的像。现在命题6由 \((\mathbf{Z}/n\mathbf{Z})^*\) 的阶为 \(\varphi(n)\) 这一事实得出。

定理 2（高斯）. --- 设 \(\bar{\mathbf{Q}}\) 为 \(\mathbf{Q}\) 的一个代数闭包，且令 \(n \geqslant 1\) 为一个整数。

\begin{enumerate}
\def\labelenumi{\alph{enumi})}
\item
  分圆多项式 \(\Phi_n(\mathbf{X})\) 在 \(\mathbf{Q}[\mathbf{X}]\) 中不可约。
\item
  \(\mathbf{R}_n(\mathbf{Q})\) 在 \(\mathbf{Q}\) 上的次数是 \(\varphi(n)\) 。
\item
  同态 \(\chi_{n}\) （从 \(\operatorname{Gal}(\bar{\mathbf{Q}} / \mathbf{Q})\) 到 \((\mathbf{Z} / n\mathbf{Z})^{*}\) 中）是满射的，并通过过渡到商群定义了 \(\operatorname{Gal}(\mathrm{R}_n(\mathbf{Q}) / \mathbf{Q})\) 到 \((\mathbf{Z} / n\mathbf{Z})^{*}\) 上的一个同构。
\end{enumerate}

考虑到命题6，我们只需证明 \(c\) ）。每个整数 \(r\) （与 \(n\) 互素）是一些素数 \(p_1, \ldots, p_s\) （均不整除 \(n\) ）的乘积；因此只需证明，对每个素数 \(p\) （不整除 \(n\) ），映射 \(x \mapsto x^p\) （从 \(\mu_n(\mathbf{Q})\) 到其自身）可延拓为 \(\mathbf{R}_n(\mathbf{Q})\) 的一个自同构。只需证明：若 \(\zeta\) 是一个本原的 \(n\) 次单位根，则极小多项式 \(f\) （ \(\zeta\) 在 \(\mathbf{Q}\) 上的）等于极小多项式 \(g\) （ \(\zeta^p\) 在 \(\mathbf{Q}\) 上的）。

我们将假设 \(f \neq g\) 并用反证法论证。多项式 \(f\) 和 \(g\) 在 \(\mathbf{Q}[\mathbf{X}]\) 中是首一不可约的，且都整除 \(\mathbf{X}^n - 1\) ，因此存在 \(u \in \mathbf{Q}[\mathbf{X}]\) 使得 \(\mathbf{X}^n - 1 = fgu\) （IV，第13页，命题13）。引理2（V，第82页）表明 \(f, g\) 和 \(u\) 具有整数系数。我们用 \(\bar{v}\) 表示系数在 \(\mathbf{F}_p\) 中的、由多项式 \(v \in \mathbf{Z}[\mathbf{X}]\) 经模 \(p\) 约化得到的多项式。于是我们有 \(\mathbf{X}^n - 1 = \bar{f}\bar{g}\bar{u}\) （在 \(\mathbf{F}_p[\mathbf{X}]\) 中）。

进一步，我们有 \(g(\zeta^p) = 0\) ，因此 \(g(\mathbf{X}^p)\) 是 \(f(\mathbf{X})\) 在 \(\mathbf{Q}[\mathbf{X}]\) 中的倍式。根据引理2，存在 \(h \in \mathbf{Z}[\mathbf{X}]\) 使得 \(g(\mathbf{X}^p) = f(\mathbf{X}) \cdot h(\mathbf{X})\) 。现在我们有 \(v(\mathbf{X}^p) = v(\mathbf{X})^p\) （对每个多项式 \(v \in \mathbf{F}_p[\mathbf{X}]\) ）。通过模 \(p\) 约化，我们得到 \(\overline{g}^p = \overline{f}\overline{h}\) 。若 \(v\) 是 \(\mathbf{F}_p[\mathbf{X}]\) 中整除 \(\overline{f}\) 的一个不可约多项式，则它必定整除 \(\overline{g}\) 。由于 \(\overline{f}\overline{g}\) 整除 \(\mathbf{X}^n - 1\) ，我们得出 \(v^2\) 整除 \(\mathbf{X}^n - 1\) （在 \(\mathbf{F}_p[\mathbf{X}]\) 中）。这是荒谬的，因为多项式 \(\mathbf{X}^n - 1\) 在 \(\mathbf{F}_p[\mathbf{X}]\) 中是可分的。

可以证明，对于 Q 上每个有限次阿贝尔扩张 E，存在一个整数 \(n \geqslant 1\) 使得 E 同构于 \(\mathrm{R}_{n}(\mathbf{Q})\) 的一个子扩张。* 换言之，域 \(\mathbf{Q}(\mu_{\infty}(\mathbf{C}))\) 是 \(\mathbf{Q}\) 的一个阿贝尔闭包。*（«克罗内克–韦伯定理»）。

